\begin{abstract}
Prediction markets provide well-calibrated probability estimates for future events, yet raw probabilities lack the narrative context needed for broad consumption. We present \systemname, an end-to-end system that combines large language models with prediction market data to generate plausible news articles about future events. Using 25 resolved \polymarket events spanning politics, sports, entertainment, and geopolitics---with a combined trading volume of \$12.6 billion---we evaluate four generation strategies: an unconditioned baseline, direct market conditioning, Market-Conditioned Prompting (\mcp), and a superforecaster persona. We find that market conditioning dramatically improves factual accuracy, with conditioned articles scoring 4.64--4.88 out of 5 compared to 1.72 for unconditioned generation (Cohen's $d = 2.3$--$2.9$, $p < 0.0001$). Without market data, \gptfour produces well-written articles about the \emph{wrong} future; with market data, it produces well-written articles about the \emph{right} future. Writing quality remains uniformly high across all conditions (5.0/5 coherence), confirming that narrative generation and factual grounding are orthogonal capabilities. Different conditioning strategies offer distinct trade-offs: direct conditioning maximizes accuracy (4.88/5), while the superforecaster persona maximizes informativeness (5.0/5). Our results suggest that forecasting and narrative generation should remain separate pipeline stages, and that prediction market data can be effectively translated into accessible news narratives.
\end{abstract}
